\section{Q2}

\subsection{Post Selected Pauli Projection}

Note we have the following identity

\begin{align*}
	\tikzfig{q-2-1.1} 
\end{align*}

So we have 
\ctikzfig{q-2-1.2}

where $\vec{Q} = Z \otimes \vec{P}$


\subsection{Two Pauli and Clifford}

The goal of this section is to show that for any $j, k \in \{0, 1\}$, we have Pauli $\vec{Q}_1, \vec{Q}_2$ and Clifford $C_{i, j}$ such that the following identity holds 

\begin{equation}\label{question:q2}
	\tikzfig{q-2-question}
\end{equation}

The desired Paulis are $\vec{Q_1} = Z \otimes \vec{P}$, $\vec{Q_2} = X \otimes \vec{I}$.
The desired Cliffords are shown in the following cases.

\subsubsection{Case 0,0}
We have
$C_{0,0} = H \otimes \vec{I}$.

Plugin to $\eqref{question:q2}$, we have
\begin{align*}
	\tikzfig{q-2-2.1}
\end{align*}
 
which is desired.

\subsubsection{Case 0,1}

I claim that $C_{0,1} = (X H) \otimes P$ works, where $P$ is the Pauli corresponding to the initial Pauli measurement.

Check the following diagram
\begin{align*}
	\tikzfig{q-2-2.2}
\end{align*}

Where in the last step we used the fact that $P = \bigotimes P_i$ and $PP = I$ (as $P$ is Pauli).
In the second last step we use the following identity

\ctikzfig{q-2-2-Pi-measure}

\begin{proof}
\begin{align*}
	\tikzfig{q-2-2.3}
\end{align*}
\end{proof}

\subsubsection{Case 1,0}

Using the same argument as above we see 
I claim that $C_{1,0} = I \otimes P$.

\subsubsection{Case 1,1}

Using the same argument as above we see 
I claim that $C_{1,1} = (X H) \otimes I$.
